\documentclass[a4paper,12pt,twoside]{memoir}

% Castellano
\usepackage[spanish,es-tabla]{babel}
\selectlanguage{spanish}
\usepackage[utf8]{inputenc}
\usepackage[T1]{fontenc}
\usepackage{lmodern} % Scalable font
\usepackage{microtype}
\usepackage{placeins}

\RequirePackage{booktabs}
\RequirePackage[table]{xcolor}
\RequirePackage{xtab}
\RequirePackage{multirow}

% Links
\PassOptionsToPackage{hyphens}{url}\usepackage[colorlinks]{hyperref}
\hypersetup{
	allcolors = {red}
}

% Ecuaciones
\usepackage{amsmath}

% Rutas de fichero / paquete
\newcommand{\ruta}[1]{{\sffamily #1}}

% Párrafos
\nonzeroparskip

% Huérfanas y viudas
\widowpenalty100000
\clubpenalty100000

% Imágenes

% Comando para insertar una imagen en un lugar concreto.
% Los parámetros son:
% 1 --> Ruta absoluta/relativa de la figura
% 2 --> Texto a pie de figura
% 3 --> Tamaño en tanto por uno relativo al ancho de página
\usepackage{graphicx}
\newcommand{\imagen}[3]{
	\begin{figure}[!h]
		\centering
		\includegraphics[width=#3\textwidth]{#1}
		\caption{#2}\label{fig:#1}
	\end{figure}
	\FloatBarrier
}

% Comando para insertar una imagen sin posición.
% Los parámetros son:
% 1 --> Ruta absoluta/relativa de la figura
% 2 --> Texto a pie de figura
% 3 --> Tamaño en tanto por uno relativo al ancho de página
\newcommand{\imagenflotante}[3]{
	\begin{figure}
		\centering
		\includegraphics[width=#3\textwidth]{#1}
		\caption{#2}\label{fig:#1}
	\end{figure}
}

% El comando \figura nos permite insertar figuras comodamente, y utilizando
% siempre el mismo formato. Los parametros son:
% 1 --> Porcentaje del ancho de página que ocupará la figura (de 0 a 1)
% 2 --> Fichero de la imagen
% 3 --> Texto a pie de imagen
% 4 --> Etiqueta (label) para referencias
% 5 --> Opciones que queramos pasarle al \includegraphics
% 6 --> Opciones de posicionamiento a pasarle a \begin{figure}
	\newcommand{\figuraConPosicion}[6]{%
		\setlength{\anchoFloat}{#1\textwidth}%
		\addtolength{\anchoFloat}{-4\fboxsep}%
		\setlength{\anchoFigura}{\anchoFloat}%
		\begin{figure}[#6]
			\begin{center}%
				\Ovalbox{%
					\begin{minipage}{\anchoFloat}%
						\begin{center}%
							\includegraphics[width=\anchoFigura,#5]{#2}%
							\caption{#3}%
							\label{#4}%
						\end{center}%
					\end{minipage}
				}%
			\end{center}%
		\end{figure}%
	}
	
	%
	% Comando para incluir imágenes en formato apaisado (sin marco).
	\newcommand{\figuraApaisadaSinMarco}[5]{%
		\begin{figure}%
			\begin{center}%
				\includegraphics[angle=90,height=#1\textheight,#5]{#2}%
				\caption{#3}%
				\label{#4}%
			\end{center}%
		\end{figure}%
	}
	% Para las tablas
	\newcommand{\otoprule}{\midrule [\heavyrulewidth]}
	%
	% Nuevo comando para tablas pequeñas (menos de una página).
	\newcommand{\tablaSmall}[5]{%
		\begin{table}
			\begin{center}
				\rowcolors {2}{gray!35}{}
				\begin{tabular}{#2}
					\toprule
					#4
					\otoprule
					#5
					\bottomrule
				\end{tabular}
				\caption{#1}
				\label{tabla:#3}
			\end{center}
		\end{table}
	}
	
	%
	% Nuevo comando para tablas pequeñas (menos de una página).
	\newcommand{\tablaSmallSinColores}[5]{%
		\begin{table}[H]
			\begin{center}
				\begin{tabular}{#2}
					\toprule
					#4
					\otoprule
					#5
					\bottomrule
				\end{tabular}
				\caption{#1}
				\label{tabla:#3}
			\end{center}
		\end{table}
	}
	
	\newcommand{\tablaApaisadaSmall}[5]{%
		\begin{landscape}
			\begin{table}
				\begin{center}
					\rowcolors {2}{gray!35}{}
					\begin{tabular}{#2}
						\toprule
						#4
						\otoprule
						#5
						\bottomrule
					\end{tabular}
					\caption{#1}
					\label{tabla:#3}
				\end{center}
			\end{table}
		\end{landscape}
	}
	
	%
	% Nuevo comando para tablas grandes con cabecera y filas alternas coloreadas en gris.
	\newcommand{\tabla}[6]{%
		\begin{center}
			\tablefirsthead{
				\toprule
				#5
				\otoprule
			}
			\tablehead{
				\multicolumn{#3}{l}{\small\sl continúa desde la página anterior}\\
				\toprule
				#5
				\otoprule
			}
			\tabletail{
				\hline
				\multicolumn{#3}{r}{\small\sl continúa en la página siguiente}\\
			}
			\tablelasttail{
				\hline
			}
			\bottomcaption{#1}
			\rowcolors {2}{gray!35}{}
			\begin{xtabular}{#2}
				#6
				\bottomrule
			\end{xtabular}
			\label{tabla:#4}
		\end{center}
	}
	
	%
	% Nuevo comando para tablas grandes con cabecera.
	\newcommand{\tablaSinColores}[6]{%
		\begin{center}
			\tablefirsthead{
				\toprule
				#5
				\otoprule
			}
			\tablehead{
				\multicolumn{#3}{l}{\small\sl continúa desde la página anterior}\\
				\toprule
				#5
				\otoprule
			}
			\tabletail{
				\hline
				\multicolumn{#3}{r}{\small\sl continúa en la página siguiente}\\
			}
			\tablelasttail{
				\hline
			}
			\bottomcaption{#1}
			\begin{xtabular}{#2}
				#6
				\bottomrule
			\end{xtabular}
			\label{tabla:#4}
		\end{center}
	}
	
	%
	% Nuevo comando para tablas grandes sin cabecera.
	\newcommand{\tablaSinCabecera}[5]{%
		\begin{center}
			\tablefirsthead{
				\toprule
			}
			\tablehead{
				\multicolumn{#3}{l}{\small\sl continúa desde la página anterior}\\
				\hline
			}
			\tabletail{
				\hline
				\multicolumn{#3}{r}{\small\sl continúa en la página siguiente}\\
			}
			\tablelasttail{
				\hline
			}
			\bottomcaption{#1}
			\begin{xtabular}{#2}
				#5
				\bottomrule
			\end{xtabular}
			\label{tabla:#4}
		\end{center}
	}
	
	
	
	\definecolor{cgoLight}{HTML}{EEEEEE}
	\definecolor{cgoExtralight}{HTML}{FFFFFF}
	
	%
	% Nuevo comando para tablas grandes sin cabecera.
	\newcommand{\tablaSinCabeceraConBandas}[5]{%
		\begin{center}
			\tablefirsthead{
				\toprule
			}
			\tablehead{
				\multicolumn{#3}{l}{\small\sl continúa desde la página anterior}\\
				\hline
			}
			\tabletail{
				\hline
				\multicolumn{#3}{r}{\small\sl continúa en la página siguiente}\\
			}
			\tablelasttail{
				\hline
			}
			\bottomcaption{#1}
			\rowcolors[]{1}{cgoExtralight}{cgoLight}
			
			\begin{xtabular}{#2}
				#5
				\bottomrule
			\end{xtabular}
			\label{tabla:#4}
		\end{center}
	}
	
	
	
	\graphicspath{ {./img/} }
	
	% Capítulos
	\chapterstyle{bianchi}
	\newcommand{\capitulo}[2]{
		\setcounter{chapter}{#1}
		\setcounter{section}{0}
		\setcounter{figure}{0}
		\setcounter{table}{0}
		\chapter*{#2}
		\addcontentsline{toc}{chapter}{#2}
		\markboth{#2}{#2}
	}
	
	% Apéndices
	\renewcommand{\appendixname}{Apéndice}
	\renewcommand*\cftappendixname{\appendixname}
	
	\newcommand{\apendice}[1]{
		%\renewcommand{\thechapter}{A}
		\chapter{#1}
	}
	
	\renewcommand*\cftappendixname{\appendixname\ }
	
	% Formato de portada
	\makeatletter
	\usepackage{xcolor}
	\newcommand{\tutor}[1]{\def\@tutor{#1}}
	\newcommand{\course}[1]{\def\@course{#1}}
	\definecolor{cpardoBox}{HTML}{E6E6FF}
	\def\maketitle{
		\null
		\thispagestyle{empty}
		% Cabecera ----------------
		\noindent\includegraphics[width=\textwidth]{cabecera}\vspace{1cm}%
		\vfill
		% Título proyecto y escudo informática ----------------
		\colorbox{cpardoBox}{%
			\begin{minipage}{.8\textwidth}
				\vspace{.5cm}\Large
				\begin{center}
					\textbf{TFG del Grado en Ingeniería Informática}\vspace{.6cm}\\
					\textbf{\LARGE\@title{}}
				\end{center}
				\vspace{.2cm}
			\end{minipage}
			
		}%
		\hfill\begin{minipage}{.20\textwidth}
			\includegraphics[width=\textwidth]{escudoInfor}
		\end{minipage}
		\vfill
		% Datos de alumno, curso y tutores ------------------
		\begin{center}%
			{%
				\noindent\LARGE
				Presentado por \@author{}\\ 
				en Universidad de Burgos --- \@date{}\\
				Tutor: \@tutor{}\\
			}%
		\end{center}%
		\null
		\cleardoublepage
	}
	\makeatother
	
	\newcommand{\nombre}{David Renedo Gil} %%% cambio de comando
	\newcommand{\nombreTFG}{GII 22.24 Aplicación de gestión de TFGs}

	% Datos de portada
	\title{\nombreTFG}
	\author{\nombre}
	\tutor{Álvar Arnaiz González y Ana Serrano Mamolar}
	\date{\today}
	\begin{document}
		
		\maketitle
		
		
		\newpage\null\thispagestyle{empty}\newpage
		
		
		%%%%%%%%%%%%%%%%%%%%%%%%%%%%%%%%%%%%%%%%%%%%%%%%%%%%%%%%%%%%%%%%%%%%%%%%%%%%%%%%%%%%%%%%
		\thispagestyle{empty}
		
		
		\noindent\includegraphics[width=\textwidth]{cabecera}\vspace{1cm}
		
		\noindent D. Álvar Arnaiz  Gonzálezy Dña. Ana Serrano Mamolar, profesores del departamento de Ingeniería Informática, área de Lenguajes y Sistemas Informáticos.
		
		\noindent Expone:
		
		\noindent Que el alumno D. \nombre, con DNI dni, ha realizado el Trabajo final de Grado en Ingeniería Informática titulado GII 22.24 Aplicación de gestión de TFGs.
		
		\noindent Y que dicho trabajo ha sido realizado por el alumno bajo la dirección del que suscribe, en virtud de lo cual se autoriza su presentación y defensa.
		
		\begin{center} %\large
			En Burgos, {\large \today}
		\end{center}
		
		\vfill\vfill\vfill
		
		% Author and supervisor
		\begin{minipage}{0.45\textwidth}
			\begin{flushleft} %\large
				Vº. Bº. del Tutor:\\[2cm]
				Dña. Ana Serrano Mamolar
			\end{flushleft}
		\end{minipage}
		\hfill
		\begin{minipage}{0.45\textwidth}
			\begin{flushleft} %\large
				Vº. Bº. del co-tutor:\\[2cm]
				D. Álvar Arnaiz González
			\end{flushleft}
		\end{minipage}
		\hfill
		
		\vfill
		
		% para casos con solo un tutor comentar lo anterior
		% y descomentar lo siguiente
		%Vº. Bº. del Tutor:\\[2cm]
		%D. nombre tutor
		
		
		\newpage\null\thispagestyle{empty}\newpage
		
		
		
		
		\frontmatter
		
		% Abstract en castellano
		\renewcommand*\abstractname{Resumen}
		\begin{abstract}
			El proyecto propone evolucionar la aplicación web de Gestor de trabajos de fin de grado, una aplicación web para el manejo de procesos de oferta, búsqueda, asignación y evaluación de los TFG de la carrera de Ingeniería Informática en la Universidad de Burgos. Usando como lenguaje de programación Java y empleando Vaadin como framework para los componentes gráficos.
			Se requieren mejoras visuales, así como la creación de nuevas pantallas junto con la implantación de \emph{websrcaping} para crear informes y presentar datos, y la posibilidad de asignar y ofertar TFGs directamente desde la aplicación. También se da la oportunidad de cambiar la información de los TFGs activos a los usuarios administradores.
		\end{abstract}
		
		\renewcommand*\abstractname{Descriptores}
		\begin{abstract}
			Aplicación web, Vaadin, Eclipse IDE, Java, trabajo de fin de grado, gestor de proyectos, Moodle, Maven, Apache Tomcat, WebScraping, framework, formularios.
		\end{abstract}
		
		\clearpage
		
		% Abstract en inglés
		\renewcommand*\abstractname{Abstract}
		\begin{abstract}
			The project proposes to evolve the fina degree prject manager web application, in charge of managing the offer, search, assignment and evaluation of Computer Engineering career final projects career of the University of Burgos. It uses Java as programming language and employs graphic components of the Vaadin framework.
			Visual improvements are required, as well as the creation of new screens along with the implementation of websrcaping to create reports and present data, and the possibility of assigning and bidding TFGs directly from the application. It is also given the opportunity to change the information of the active projects by the administrator users.
		\end{abstract}
		
		\renewcommand*\abstractname{Keywords}
		\begin{abstract}
			Web Application, Vaadin, Eclipse IDE, Java, final degree project, project manager, Moodle, Maven, Apache Tomcat, WebScraping, framework, forms.
		\end{abstract}
		
		\clearpage
		
		% Indices
		\tableofcontents
		
		\clearpage
		
		\listoffigures
		
		\clearpage
		
		\listoftables
		\clearpage
		
		\mainmatter
		\capitulo{1}{Introducción}

El proyecto se centrará en la mejora de la actual aplicación del Gestor de asignaciones de TFG, utilizada en el grado de Ingeniería Informática. La aplicación actual se trata de una actualización previa del TFG denominado \href{https://github.com/dbo1001/Gestor-TFG-2021}{GII 20.09 Herramienta web
repositorios de TFGII}. 



\section{Estructura de la memoria}
La memoria consta de los siguientes apartados:

\begin{itemize}
	\item \textbf{Introducción:} presentación del proyecto y estructuración de este.
	\item \textbf{Objetivos del proyecto:}  exposición de los  objetivos generales, técnicos y personales del proyecto.
	\item \textbf{Conceptos teóricos:} explicación de los términos teóricos necesarios para la comprensión y el desarrollo del proyecto.
	\item \textbf{Técnicas y herramientas:} definición de las técnicas utilizadas para el desarrollo del proyecto y de las herramientas empleadas para su funcionamiento.
	\item \textbf{Aspectos relevantes del desarrollo:} breve explicación de los términos más importantes durante el desarrollo del proyecto.
	\item \textbf{Trabajos relacionados:} descripción de los trabajos y proyectos asociados con la gestión de trabajos de fin de grado (TFG).
	\item \textbf{Conclusiones y líneas de trabajo futuras:} conclusiones obtenidas al finalizar el proyecto y descripción de posibles futuras líneas de trabajo o mejoras.
	\item \textbf{Anexo}: se realiza de forma complementaria un documento de Anexo con explicación más detalla de todo el desarrollo del proyecto.
\end{itemize}
		\capitulo{2}{Objetivos del proyecto}

Este apartado explica de forma precisa y concisa cuáles son los objetivos que se persiguen con la realización del proyecto.

\section{Objetivos generales}
\begin{itemize}
	\item Corregir diversos \emph{bugs} de la versión anterior, entre ellos la modificación de campos de la base de datos que a pesar de realizar bien los cálculos de las estadísticas, se recogían mal, así como corregir la asignación de un proyecto a un determinado curso, que en la pantalla de histórico de proyectos se realizaba de forma errónea.
	\item Añadir nuevas funcionalidades para completar la aplicación.
		\begin{itemize}
			\item Pantalla con la visualización del histórico de los profesores, áreas y departamentos.
			\item Permitir crear un informe de dedicación de datos de profesores y áreas.
			\item Implementar la opción de oferta/asignación de TFGs.
			\item Implementar la funcionalidad de aceptar/denegar TFGs propuestos, así como de modificar los proyectos de la pestaña de activos.
			\item Permitir la descarga de la base de datos utilizada.
			
		\end{itemize}

	
\end{itemize}
\section{Objetivos técnicos}
\begin{itemize}
	\item Implementar \emph{Web Scraping} con la página \href{https://investigacion.ubu.es/unidades/2682/investigadores}{de investigadores de la UBU}.
	\item Permitir la oferta de proyectos desde la propia aplicación siendo profesor y la aceptación de los mismos por parte de un administrador.
	\item Crear una nueva gráfica para las estadísticas de los profesores con las opciones marcadas por el usuario.
	\item Guardar datos generados en archivos \emph{.csv} y \emph{.xls}.
	\item Actualizar los ficheros de datos del profesorado.
	\item Utilizar GitHub para llevar a cabo el seguimiento del proyecto y control de versiones.
	\item Registrar el usuario que inicia sesión en el sistema para asignarle un rol y mostrarle las pantallas correspondientes a dicho actor, creando tres de ellos: alumno, profesor y administrador.
\end{itemize}
		\capitulo{3}{Conceptos teóricos}

En este apartado se definirán algunos de los conceptos utilizados a lo largo del proyecto.


\section{\textit{Webscraping}}
El textbf{\textit{WebScraping}}~\cite{webScrap} es una técnica que se utiliza para la extracción y almacenamiento de información de cualquier página web a través de un programa de software, que suele ser un crawler. Aunque, a priori, puede parecer que cualquiera puede scrapear información en cualquier sitio web, lo cierto es que la extracción de datos no siempre es legal. Por ejemplo, los datos que requieren un registro del usuario no pueden ser obtenidos a través del web scraping.

\section{\textit{Desarrollo Backend}}
Dentro del desarrollo web, el \textbf{\textit{backend}}~\cite{backend} se encarga de todos los procesos necesarios para que la web funcione de forma correcta. Estos procesos o funciones no son visibles, pero tienen mucha importancia en el buen funcionamiento de un sitio web. Algunas de estas acciones que controla el backend son la conexión con la base de datos o la comunicación con el servidor de hosting.

\section{\textit{Desarrollo Frontend}}
\textbf{\textit{Frontend}}~\cite{Frontend} es la parte de una aplicación que interactúa con los usuarios, es conocida como el lado del cliente. Básicamente es todo lo que vemos en la pantalla cuando accedemos a un sitio web o aplicación: tipos de letra, colores, adaptación para distintas pantallas(RWD), los efectos del ratón, teclado, movimientos, desplazamientos, efectos visuales… y otros elementos que permiten navegar dentro de una página web. Este conjunto crea la experiencia del usuario.

\section{\textit{Framework}}
Los \textbf{\textit{framework}}~\cite{frameworks} web son un conjunto de herramientas, estilos y librerías dispuestas a través de una estructura o esqueleto base, para el desarrollo de aplicaciones web más escalables y sencillas de mantener.
Gracias a estos frameworks web, podemos ahorrar grandes cantidades de tiempo y costes, pero vamos a profundizar más en las ventajas que tienen, causantes de su gran éxito y expansión.




		\capitulo{4}{Técnicas y herramientas}
En esta sección se detallarán algunas de las técnicas y herramientas utilizadas durante el desarrollo del trabajo y las alternativas que no se han llegado a usar.

Muchas de las herramientas utilizadas están ya explicadas en la anterior versión del proyecto que podemos encontrar en \href{https://github.com/dbo1001/Gestor-TFG-2021}{GII 20.09 Herramienta web repositorios de TFGII}. Por lo tanto lenguajes de programación como \emph{Java, SQL, XML} no se detallarán a pesar de haber sido utilizados, y herramientas como \emph{GitHub, GitHub Desktop, Eclipse IDE, Tomcat o Maven} tampoco se explicarán.

\section{Herramientas de Desarrollo}

\subsection{Heroku para estudiantes}
Heroku~\cite{heroku_students} es una plataforma basada en la nube como servicio ( PaaS ) para construir, ejecutar y administrar aplicaciones. 
La plataforma Integración de GitHub le permite conectar su aplicación Heroku a su repositorio GitHub e implementar en cada empuje a GitHub.
Heroku ofrece una gama de servicios de bajo costo para ayudarlo a experimentar, aprender y crear prototipos de nuevas ideas. Para los estudiantes de GitHub, vamos un paso más allá y agregamos aún más recursos a su Paquete de desarrolladores de estudiantes de GitHub. 
Heroku pasó a ser de pago el 28 de noviembre de 2022,pero al pertenecer al programa de GitHub para estudiantes proporcionado por la UBU, podemos solicitarlo, y que se nos agreguen unos créditos para poder desplegar nuestro proyecto en la nube,por lo que es la herramienta que hemos utilizado para desplegar nuestro proyecto.

\subsection{Jsoup}
Jsoup~\cite{jsoup} es una biblioteca Java para trabajar con HTML del mundo real. Proporciona una API muy conveniente para recuperar URL y extraer y manipular datos, utilizando los mejores métodos HTML5 DOM y selectores CSS. 
Es la herramienta que hemos utilizado para realizar el \emph{webscraping} ya que es la que mejor se ajustaba y con la que mejor resultado hemos obtenido, también resulta la más sencilla de utilizar.

\subsection{HTMLUnit}
HtmlUnit~\cite{HTMLUNIT} es un "navegador sin GUI para programas Java". Modela documentos HTML y proporciona una API que le permite invocar páginas, completar formularios, hacer clic en enlaces, etc... tal como lo haces en tu navegador "normal. 
Es una de las alternativas planteadas a la hora de utilizar una biblioteca para el \emph{webscrap}.

\subsection{Pencil}
Pencil~\cite{Pencil_project} está construido con el propósito de proporcionar una herramienta de creación de prototipos GUI gratuita y de código abierto que las personas pueden instalar y usar fácilmente para crear maquetas en plataformas de escritorio populares. 
Se ha utilizado para crear los \emph{mock-ups} de las nuevas pantallas que se han desarrollado.

\subsection{Northflank}
Northflank~\cite{Northflank} es su sitio web cabe mencionar que los precios escalan junto a como va desarrollándose tu proyecto. Puedes pagar solo por los recursos que consumen tus servicios durante la compilación y el tiempo de ejecución. Puedes escalar tanto horizontal como verticalmente sin gastos inesperados. Al ser un VPS, puedes trabajar con muchos lenguajes y frameworks. Cuentan con una versión para Developers que es gratuita. 
Finalmente se optó por no utilizar este sistema ya que tras intentar importar nuestro proyecto nos indicaba que era de pago por que superaba el límite de peso permitido para un proyecto gratuito.

\subsection{Vaadin}
Vaadin~\cite{Vaadin} es un framework de desarrollo de SPA que permite escribir el código de dichas aplicaciones en Java o en cualquier otro lenguaje soportado por la JVM 1.6+. Esto permite la programación de la interfaz gráfica en lenguajes como Java 8, Scala o Groovy, por ejemplo.

Uno de las características diferenciadores de Vaadin es que, contrario a las librerías y frameworks de JavaScript típicas, presenta una arquitectura centrada en el servidor, lo que implica que la mayoría de la lógica es ejecutada en los servidores remotos. Del lado del cliente, Vaadin está construido encima de Google Web Toolkit, con el que puede extenderse.

\subsection{LucidChart}

LucidChart ~\cite{LucidChart} es una aplicación de diagramación inteligente que reúne a los equipos para que tomen mejores decisiones y construyan el futuro. Ha sido utilizada para realizar los diagramas de flujo del Anexo.

\section{Herramientas de Documentación}

\subsection{Latex}
 Latex~\cite{latex} es un sistema de composición de textos, orientado a la creación de documentos escritos que presenten una alta calidad tipográfica. Por sus características y posibilidades, es usado de forma especialmente intensa en la generación de artículos y libros científicos que incluyen, entre otros elementos, expresiones matemáticas.
 
\subsection{TexStudio}
TexStudio~\cite{TexStudio} es un entorno de escritura integrado para crear documentos LaTeX. Su principal objetivo es hacer que escribir LaTeX sea lo más fácil y cómodo posible. Por esto, TeXstudio nos ofrece a los usuarios numerosas funciones como el resaltado de sintaxis, un visor integrado, la verificación de referencias y varios asistentes, entre otras. 
Es el programa de edición que hemos utilizado para realizar la memoria y los anexos.

\subsection{MikTex}
MiKTeX~\cite{TexStudio} es una implementación de TeX y programas relacionados para Windows. TeX es un sistema de caracteres de alta calidad para la edición de documentos.
Es la herramienta que se utiliza para compilar los documentos que se crean en TexStudio.

\subsection{Overleaf}
Overleaf~\cite{Overleaf}  es una herramienta de publicación y redacción colaborativa en línea que hace que todo el proceso de redacción, edición y publicación de documentos científicos sea mucho más rápido y sencillo. Overleaf brinda la conveniencia de un editor LaTeX fácil de usar con colaboración en tiempo real y la salida totalmente compilada producida automáticamente en segundo plano a medida que escribe.
Era la herramienta que se iba a utilizar en un principio pero se deshechó la idea ya que se vió que era más intuitivo de utilizar la aplicación de escritorio TexStudio.


		\capitulo{5}{Aspectos relevantes del desarrollo del proyecto}

Este apartado pretende recoger los aspectos más interesantes del desarrollo del proyecto, comentados por los autores del mismo.
Debe incluir desde la exposición del ciclo de vida utilizado, hasta los detalles de mayor relevancia de las fases de análisis, diseño e implementación.
Se busca que no sea una mera operación de copiar y pegar diagramas y extractos del código fuente, sino que realmente se justifiquen los caminos de solución que se han tomado, especialmente aquellos que no sean triviales.
Puede ser el lugar más adecuado para documentar los aspectos más interesantes del diseño y de la implementación, con un mayor hincapié en aspectos tales como el tipo de arquitectura elegido, los índices de las tablas de la base de datos, normalización y desnormalización, distribución en ficheros3, reglas de negocio dentro de las bases de datos (EDVHV GH GDWRV DFWLYDV), aspectos de desarrollo relacionados con el WWW...
Este apartado, debe convertirse en el resumen de la experiencia práctica del proyecto, y por sí mismo justifica que la memoria se convierta en un documento útil, fuente de referencia para los autores, los tutores y futuros alumnos.

		\capitulo{6}{Trabajos relacionados}
Se nombrarán algunos de los proyectos y aplicaciones similares o relacionados con la gestión de trabajos de fin de grado o master que se han visitado para la obtención de información del proyecto.

Varios de los TFGs visualizados son los mismos que menciona Diana en la \href{https://github.com/dbo1001/Gestor-TFG-2021}{versión anterior},de la que se hablará, por lo que no se comentarán.  

\section{GII 20.09 Herramienta web repositorios de TFGII}
Este el proyecto que se propone para mejorar consiste en una aplicación web para la gestión de TFG de la Universidad de Burgos.

Este proyecto es una actualización de la versión previa encontrada en \href{https://github.com/jfb0019/Gestor-TFG-2016}{repositorio de GitHub}. 

El código fuente se puede encontrar en este \href{https://github.com/dbo1001/Gestor-TFG-2021}{repositorio de GitHub}. En el anexo se ha explicado como lo hemos compilado y ejecutado para ver su funcionamiento y seguir con el desarrollo. Es un trabajo realizado mediante Java 11, con Vaadin 14 como plataforma de código abierto para la interfaz web. 

\section{title}

		\capitulo{7}{Conclusiones y Líneas de trabajo futuras}

\section{Conclusiones}

Una vez finalizado el proyecto y remontando la vista atrás al comienzo del mismo puedo afirmar que se han cumplido todos los objetivos planificados y, además, obtenido muchos conocimientos.

La mayoría de ellos están relacionados con el \emph{framework de Vaadin} y su funcionamiento.
El proyecto ha cumplido con todas las propuestas planteadas al inicio del curso y, además, se han ido supervisando las diferentes formas de adaptarlas mediante ideas de los tutores y mías propias. 
Por ejemplo, en un primer momento sólo se planteaba la opción de proponer proyectos y posteriormente darle a un administrador la oportunidad de aceptarlos y denegarlos. Pero en la etapa final del mismo también se decidió replantearlo y permitir al administrador modificar los proyectos que estuviesen activos que quisiese, de esta forma se ha tenido que añadir una funcionalidad extra al proceso.

El proyecto propone mejorar una aplicación ya existente por lo que gran parte del comienzo implica estudiar y comprender cómo funciona y cómo está estructurada la versión anterior. La mayor parte del estudio se ve facilitada gracias a los comentarios que estaban expuestos en el código y que también se han aumentado en esta versión, como ayuda para posteriores mejoras.
Se empleó mucho tiempo en conocer el código y sobretodo en entender las herramientas de uso, entre ellas destacar el uso y funcionamiento del ya mencionado \emph{Vaadin} y de la propia herramienta de documentación \LaTeX{} que nunca había sido utilizada.

Gran parte del funcionamiento se basó en la extracción de información sobre el profesorado de la Escuela Politécnica Superior mediante \emph{webscraping}, cuyos datos se utilizan en la mayoría de las pantallas añadidas. Se estudió cómo realizar esta técnica, y se intentó mediante diferentes librerías dando finalmente con la solución.

He aprendido mucho sobre el funcionamiento de la aplicación y de los protocolos utilizados para realizar las peticiones entre cliente y servidor. En cuanto a los lenguajes, he reforzado los conocimientos que tenía sobre SQL y sobretodo Java, que se cursan a lo largo del grado. También he aprendido un poco sobre CSS, del cual tenía unos conocimientos muy vagos.

Otro de los aspectos más relevantes que se ha tenido que tratar ha sido la gestión del tipo de usuario que iniciaba sesión en la aplicación, cómo obtener su rol y de esta forma restringir las vistas a las que puede acceder. Se estudiaron numerosas alternativas para lograr esta funcionalidad, como la que se puede observar en la rama \emph{TestBranch} del proyecto y finalmente se implantó la más sencilla y efectiva.

\section{Líneas de trabajo futuras}

Algunas de las posibles líneas de trabajo futuras del programa son:

\begin{itemize}
	\item Incluir la \textbf{internacionalización} de la aplicación, traducirlo a otros idiomas como inglés.
	\item Añadir una funcionalidad que \emph{avisase a los tutores de un TFG sobre su estado}. Es decir, enviar un mensaje al tutor de un proyecto que haya sido modificado por un administrador, ya sea modificando su estado entre aceptado y denegado o cambiando su información.
	\item Con los datos obtenidos de las Cookies se podría obtener la región
	desde la cual se está accediendo y aplicarlo para calcular la fecha de
	la última modificación de los ficheros, según la zona horaria.
	\item Mejorar los aspectos visuales de las nuevas pantallas añadidas, ya que en esta versión se ha tenido más en cuenta la funcionalidad.
	\item Permitir a un alumno solicitar un proyecto de la lista de activos que existan desde la propia aplicación, que no lo tenga que hacer mediante un correo referenciando dicho proyecto. Es decir añadir una columna en el apartado de activos que pueda ser <<Solicitado por>>. De esta forma podemos diferenciar un usuario que inicia sesión y tiene rol de alumno de un usuario que no inicia sesión, que no podría solicitar proyectos.
	\item Modificar la operación de actualización de la base de datos del profesorado, \emph{webscraping}, y hacerla asíncrona, de tal manera que la aplicación no pare de funcionar mientras se realiza esta operación.
	
\end{itemize}
		
		
		\bibliographystyle{plain}
		\bibliography{bibliografia}
		
	\end{document}